Un alambre ocupa el cuarto de circunferencia $x^2+y^2=R^2$ en el primer cuadrante, donde $R>0$ es una longitud. Su densidad lineal es
\[
\rho(x,y)=\rho_0\left(1+\frac{x}{R}\right),\qquad \rho_0>0,
\]
con $\rho_0$ en unidades de masa por longitud. Calcule su masa y su densidad media $\rho_{\mathrm{med}}=m/L(C)$. Indique en qué extremo es mayor la densidad.
